\section{Conclusion and Future Work}
\label{sec:conclusion}

\subsection{Conclusion}
\label{sec:concl_summary}

This paper presented a two-phase framework for brain tumor MRI classification based
on separability-guided handcrafted feature optimization and comprehensive machine
learning classifier benchmarking. The proposed methodology was designed to address
a key limitation of existing handcrafted feature-based approaches, namely the
dependence of feature optimization on classifier performance. By separating feature
configuration from classifier evaluation, the framework enables a more objective
assessment of feature quality and provides a clearer understanding of the
relationship between feature representation and classification performance.

In the first phase, a classifier-independent optimization strategy was introduced
to identify the most discriminative configurations of four widely used handcrafted
descriptor families: Gray-Level Co-occurrence Matrix (GLCM), Local Binary Patterns
(LBP), Discrete Wavelet Transform (DWT), and Histogram of Oriented Gradients (HOG).
Candidate configurations were evaluated using a composite separability score derived
from Fisher Discriminant Ratio, Mutual Information, and Davies--Bouldin Index. The
results demonstrated that descriptor parameterization significantly affects feature
discriminability and that systematic optimization can substantially improve
representation quality compared with default settings.

In the second phase, the optimized descriptors were employed to construct multiple
feature representations and evaluate eight machine learning classifiers, including
Support Vector Machine, K-Nearest Neighbors, Random Forest, XGBoost, LightGBM,
Logistic Regression, Multi-Layer Perceptron, and Extra Trees. The benchmarking
study revealed that optimized handcrafted features provide strong discriminative
capability across diverse learning paradigms and that feature fusion strategies
generally outperform individual descriptor families by combining complementary
image characteristics.

The experimental findings highlighted three principal observations. First, feature
optimization plays a crucial role in maximizing classification performance and should
be considered an integral component of the learning pipeline rather than a fixed
preprocessing step. Second, texture-based descriptors emerged as particularly
informative for brain tumor MRI classification, while frequency and gradient
features contributed complementary information that enhanced hybrid representations.
Third, ensemble learning methods consistently demonstrated strong predictive
performance when combined with optimized handcrafted features, confirming their
suitability for high-dimensional medical image analysis tasks.

To ensure the reliability of the reported results, a comprehensive statistical
validation framework was incorporated. Cohen's $\kappa$ coefficients, bootstrap
confidence intervals, McNemar's tests, Wilcoxon signed-rank tests, and
Friedman--Nemenyi analyses collectively confirmed the robustness and statistical
significance of the observed performance differences. These analyses provide strong
evidence that the reported improvements are meaningful and reproducible rather than
resulting from random experimental variation.

Overall, the proposed framework demonstrates that carefully optimized handcrafted
feature engineering remains a highly competitive approach for brain tumor MRI
classification. Beyond achieving strong predictive performance, the methodology
offers important advantages in terms of interpretability, computational efficiency,
reproducibility, and suitability for limited-data environments. These characteristics
make the proposed approach particularly attractive for practical computer-aided
diagnosis systems where transparency and resource efficiency are important
considerations.

\subsection{Future Work}
\label{sec:concl_future}

Although the proposed framework achieved promising results, several opportunities
exist for further research and development.

A natural extension of this work involves the integration of deep feature
representations extracted from pretrained convolutional neural networks such as
VGG16, ResNet50, EfficientNet, and DenseNet. Comparing separability-guided
handcrafted features with deep representations under a unified benchmarking
framework would provide valuable insights into the relative strengths and
limitations of both approaches.

Another promising direction is the development of hybrid feature representations
that combine optimized handcrafted descriptors with deep features. Such systems may
benefit from the interpretability of handcrafted features while simultaneously
exploiting the powerful representation-learning capabilities of deep neural
networks.

Recent advances in Vision-Language Models (VLMs) also open new research
opportunities. Medical adaptations of multimodal architectures, including CLIP,
BioMedCLIP, and MedCLIP, could be investigated for brain tumor classification and
diagnosis support. Future studies may explore whether separability-guided
optimization principles can be extended to latent representations generated by
these models.

From a clinical perspective, future work should include evaluation on larger and
more diverse multi-center datasets acquired using different MRI scanners and imaging
protocols. Such experiments would provide stronger evidence regarding the
robustness and generalizability of the proposed methodology in real-world clinical
environments.

Finally, the proposed optimization framework could be extended beyond brain tumor
classification to other medical imaging applications, including neurological
disorders, breast cancer diagnosis, lung disease classification, and
histopathological image analysis. The generality of the separability-guided
optimization strategy suggests that it may constitute a broadly applicable
methodology for handcrafted feature engineering in medical image analysis.

In conclusion, this work establishes a systematic and statistically rigorous
framework for handcrafted feature optimization and classifier benchmarking in brain
tumor MRI classification. The findings demonstrate that feature separability
constitutes an effective optimization objective and provide a strong foundation
for future research at the intersection of handcrafted feature engineering, machine
learning, deep learning, and multimodal medical artificial intelligence.
